% !TEX program = lualatex
\documentclass{article}

\usepackage{indentfirst}
\usepackage{amssymb}
\usepackage{amsmath}

\title{Exercises from chapter 4}

\author{Kainã Alves Tureso - 15466391}

\begin{document}
    \maketitle

    \section*{Exercise 4}
    Let's define first $H^1_{D0}(\Omega) = \{v \in H^1(\Omega) : v\restriction_{\Gamma_D} = 0\}$. Take $v \in H^1_{D0}$. Multiplying it in both sides and integrating, as seen previously, we obtain:
    \begin{equation}
        \label{eq:VF1}
        \int_{\Omega} \mu \nabla u  \cdot \nabla v dx - \int_{\Gamma_D \cup \Gamma_R} \mu \nabla u \cdot \check{n} v ds = \int_\Omega fv dx
    \end{equation}
    Analysing the bundary term, noticing that $v =0$ at $\Gamma_D$:
    \[
        \int_{\Gamma_D \cup \Gamma_R} \mu \nabla u \cdot \check{n} v ds =  \int_{\Gamma_R} \mu \nabla u \cdot \check{n} v ds = \int_{\Gamma_R} (h-\gamma u)v ds
    \]
    Using a lift in $u$, i.e. $u = u_0 + u_g$ with $u_0 \in H^1_{D0}$ and $u_g \restriction_{\Gamma_D} = g$, we can rewrite \ref{eq:VF1} as 
    \[
        \int_\Omega \mu \nabla (u_0 + u_g) \cdot \nabla v dx - \int_{\Gamma_R} (h - \gamma(u_0 + u_g)) v ds = \int_\Omega fvdx 
    \]

    So we have the variational formulation:

    Find $u_0 \in H^1_{D0}$ such that $\forall v \in H^1_{D0}$:
    \[
        \int_\Omega \mu \nabla u_0 \cdot \nabla v dx + \int_{\Gamma_R} \gamma u_0 v ds = \int_\Omega fvdx - \int_\Omega \mu \nabla u_g \cdot \nabla v dx + \int_{\Gamma_R} (h - \gamma u_g) v ds
    \]
\end{document}